

In this work, we consider multi-modal \glspl{VAE} \citep{mvaeWu2018, mmvaeShi2019, sutter2021generalized, palumbo2023mmvae} as the standard modeling approach to tackle both joint and conditional generation of multiple modalities. Our goal here is to motivate the need to go beyond such a standard approach, to overcome limitations that affect multi-modal \glspl{VAE}, which result in a trade-off between generation quality and generative coherence \citep{Daunhawer2021, palumbo2023mmvae}.

Consider the random variable $X=\{X^{1},\dots,X^{M}\} \sim p_{D}(x^1,\dots,x^M)$, consisting in the set of $M$ of modalities sampled from the (unknown) multi-modal data distribution $p_{D}$. We indicate the marginal distribution of a single modality by $X^i \sim p_D^i(x^i)$ and the collection of a generic subset of modalities by $X^{A}\sim p_D^A(x^A)$, with $X^{A} \defeq \{X^i \}_{i\in A}$, where $A\subset \{1,\dots,M\}$ is a set of indexes. For example: given $A=\{1,3,5\}$, then $X^{A} = \{X^1, X^3, X^5\}$.

We begin by considering uni-modal \glspl{VAE} as particular instances of the Markov chain $X\rightarrow Z\rightarrow \hat{X}$, where $Z$ is a latent variable and $\hat{X}$ is the generated variable. Models are specified by the two conditional distributions, called the encoder $Z\g_{X=x}\sim q_\psi(z\g x)$, and the decoder $\hat{X}\g_{Z=z}\sim p_\theta(\hat{x}\g z)$. 
Given a prior distribution $p_n(z)$, the objective is to define a generative model whose samples are distributed as closely as possible to the original data. 

In the case of multi-modal \glspl{VAE}, we consider the general family of \gls{MOPOE} \citep{sutter2021generalized}, which includes as particular cases many existing variants, such as \gls{MVAE} \citep{mvaeWu2018} and \gls{MMVAE} \citep{mmvaeShi2019}. Formally,  a collection of $K$ arbitrary subsets of modalities $S=\{A_1,\dots A_K\}$, along with weighting coefficients $\omega_i\geq 0,\sum_{i=1}^{K}\omega_i=1$, define the posterior $q_\psi(z\g x)=\sum_i \omega_i q^i_{\psi^{A_i}}(z\g x^{A_i})$, with $\psi=\{\psi^1,\dots,\psi^K\}$. 
To lighten the notation, we use $q_{\psi^{A_i}}$ in place of $q^i_{\psi^{A_i}}$ noting that the various $q^i_{\psi^{A_i}}$ can have both different parameters $\psi^{A_i}$ and functional form. For example, in the \gls{MOPOE} \citep{sutter2021generalized} parametrization, we have: $q_{\psi^{A_i}}(z\g x^{A_i})=\prod_{j\in A_i}q_{\psi^j}(z\g x^{j})$. Our exposition is more general and not limited to this assumption.
The selection of the posterior can be understood as the result induced by the two step procedure where i) each subset of modalities $A_i$ is encoded into specific latent variables $Y_i\sim q_{\psi^{A_i}}(\cdot \g x^{A_i})$ and ii) the latent variable $Z$ is obtained as $Z=Y_i$ with probability $\omega_i$. 
Optimization is performed w.r.t. the following \gls{ELBO} \citep{Daunhawer2021, sutter2021generalized}:
\begin{equation}\label{eq:elbo_M}
    \mathcal{L}=\sum\limits_i \omega_i \int p_D(x)q_{\psi^{A_i}}(z\g x^{A_i})\log p_{\theta}(x|z)-\log\frac{q_{\psi^{A_i}}(z\g x^{A_i})}{p_n(z)}\dd z\dd x.
\end{equation}

%From the literature on uni-modal \glspl{VAE}, a well-known limitation, called the latent collapse problem \citep{alemi2018fixing,dieng2019avoiding}, affects the quality of latent variables $Z$. For example, consider an hypothetical case in which encoders and decoders are arbitrarily flexible: then, posteriors with zero mutual information with respect to model inputs are valid maximizers of \Cref{eq:elbo_M}. To prove this, it is sufficient to substitute the posteriors $q_{\psi^{A_i}}(z\g x^{A_i})=p_n(z)$ and $p_\theta(x|z)=p_D(x)$ into the \Cref{eq:elbo_M} to observe that the optimal value $\mathcal{L}=\int p_D(x)\log p_D(x)\dd x$ is achieved \citep{alemi2018fixing,dieng2019avoiding}. 

%The problem of loss of information is further exacerbated in the case of multi-modal \glspl{VAE}, as discussed in \citep{Daunhawer2021}. Intuitively, even if the encoders $q_{\psi^{A_i}}(z\g x^{A_i})$ carry relevant information about their inputs $X^{A_i}$, step ii) of the multi-modal encoding procedure described above induces a further information bottleneck. A fraction $\omega_i$ of the time, the latent variable $Z$ will be a copy of $Y_i$, that only provides information about the subset $X^{A_i}$. No matter how good the encoding step i) discussed above is, the information about $X^{\{1,\dots,M\}\setminus A}$ that is not contained in $X^{A_i}$ cannot be retrieved. 

A well-known limitation called the latent collapse problem \citep{alemi2018fixing,dieng2019avoiding} affects the quality of latent variables $Z$. Consider the hypothetical case of arbitrary flexible encoders and decoders: then, posteriors with zero mutual information with respect to model inputs are valid maximizers of \Cref{eq:elbo_M}. To prove this, it is sufficient to substitute the posteriors $q_{\psi^{A_i}}(z\g x^{A_i})=p_n(z)$ and $p_\theta(x|z)=p_D(x)$ into the \Cref{eq:elbo_M} to observe that the optimal value $\mathcal{L}=\int p_D(x)\log p_D(x)\dd x$ is achieved \citep{alemi2018fixing,dieng2019avoiding}.
The problem of information loss is exacerbated in the case of multi-modal \glspl{VAE} \citep{Daunhawer2021}. 
Intuitively, even if the encoders $q_{\psi^{A_i}}(z\g x^{A_i})$ carry relevant information about their inputs $X^{A_i}$, step ii) of the multi-modal encoding procedure described above induces a further information bottleneck. A fraction $\omega_i$ of the time, the latent variable $Z$ will be a copy of $Y_i$, that only provides information about the subset $X^{A_i}$. No matter how good the encoding step is, the information about $X^{\{1,\dots,M\}\setminus A}$ that is not contained in $X^{A_i}$ cannot be retrieved.

Furthermore, if the latent variable carries zero mutual information w.r.t. the multi-modal input, a coherent \textit{conditional} generation of a set of modalities given others is impossible, since $\hat{X}^{A_1}\perp X^{A_2}$ for any generic sets $A_1,A_2$. While the factorization $p_\theta(x\g z) = \prod_{i=1}^M p_{\theta^i}(x^i\g z)$, $\theta=\{\theta_1,\dots,\theta_M\}$ --- where we use $p_{\theta^i}$ instead of $p^i_{\theta^i}$ to unclutter the notation --- could enforce preservation of information and guarantee a better quality of the \textit{jointly} generated data, in practice, the latent collapse phenomenon induces multi-modal \glspl{VAE} to converge toward sub-optimal operating regime. When the posterior $q_\psi(z\g x)$ collapses onto the uninformative prior $p_n(z)$, the \gls{ELBO} in \Cref{eq:elbo_M} reduces to the sum of modality independent reconstruction terms 
$\sum\limits_i \omega_i \sum_{j\in A_i}\int p_D^{j}(x^{j})p_n(z)\left(\log p_{\theta^j}(x^j|z)\right)\dd z\dd x^{j}$. In this case, flexible decoders can similarly ignore the latent variable and converge to the solution $p_{\theta^j}(x^j|z)=p_D^{j}(x^{j})$ where, paradoxically, the quality of the approximation of the various marginal distributions is extremely high, while there is a complete lack of joint coherence. 

General principles to avoid latent collapse consist in explicitly forcing the learning of informative encoders $q_\theta(z\g x)$ via $\beta-$annealing of the \gls{KL} term in the \gls{ELBO} and the reduction of the representational power of encoders and decoders. While $\beta-$annealing has been explored in the literature \citep{mvaeWu2018} with limited improvements, reducing the flexibility of encoders/decoders clearly impacts the generation quality.
Hence the presence of a trade-off: to improve coherence, the flexibility of encoders/decoders should be constrained, which in turns hurt generative quality. This trade-off has been recently addressed in the literature of multi-modal \glspl{VAE} \citep{Daunhawer2021, palumbo2023mmvae}, but our experimental results in \Cref{sec:experiments} indicate that there is ample room for improvement, and that a new approach is truly needed.

